\section{Executable reductions and charged resources}\label{sec:reductions}
\subsection{Finite-state implementations}
A synchronized causal machine between experiments has an internal state
$r_t$. At a target step it receives the target action $a_t'$, chooses a
source action according to a parameter-independent kernel depending on
$(r_t,a_t')$, receives the source report, and outputs a target report while
updating $r_t$. Random seeds are independent of the parameter and physical
preparation. A stored seed that encodes previous reports is part of $r_t$.
Every raw source step and every retained state is charged.

This synchronized model suffices for the results below. A macro-step that
uses a random number of source steps requires an explicit causal stopping
protocol and its preparation certificate. Such a protocol can be included,
but a bound for one synchronized source step is not silently reused for it.

A target controller with state $z_t$ receives the target reports and updates
$z_t$. Its joint implementation with the machine stores $(r_t,z_t)$. Thus
$|r_t|\le R_t$ and $|z_t|\le M_t$ give the direct product bound $R_tM_t$.
This is an upper bound for one realization, not a claim that no smaller
implementation is possible.

\begin{definition}[Causal reduction certificate]\label{def:reduction}
For experiments over a common parameter space, a certificate
$\Gamma:\cE\to\cF$ consists of a parameter-independent causal machine,
a lift $\mathsf L_\Gamma$ of every admitted target controller to its source
implementation, and a monotone resource map $\rho_\Gamma$. Under target
controller $\alpha$, write $\widehat P_{\lambda}^{\Gamma,\alpha}$ for
the simulated target-history law obtained from the joint machine/source
execution. An $\varepsilon$-certificate satisfies
\begin{equation}\label{eq:simulation-error}
 \sup_{\lambda,\alpha}
 \norm{\widehat P_{\lambda}^{\Gamma,\alpha}
       -P_{\cF,\lambda}^{\alpha}}_{\TV}\le\varepsilon,
\end{equation}
and implementing a target budget $B$ costs at most $\rho_\Gamma(B)$.
The default certificate is pathwise. Expected-cost and tail-budget
certificates are different types and must be stated separately.
\end{definition}

For randomized strategy lifts, the right object is the joint causal
execution. It need not be representable by a terminal kernel on the raw
source history alone that is independent of the controller. Once the lift
and controller are fixed, its visible pushforward is unambiguous. This
avoids a misleading notation that would hide feedback inside a supposedly
static garbling.

An exact certificate has $\varepsilon=0$. The identity machine gives the
identity reduction. A reparametrization $r:\Lambda'\to\Lambda$ pulls back
an experiment family and is kept separate from data reduction. If the
reparametrization is not injective it introduces parameter redundancy;
it is not loss of observed data.

\subsection{The path-kernel estimate}
We use the convention
$\norm{P-Q}_{\TV}=\sup_B|P(B)-Q(B)|$. For $0\le f\le L$ this implies
$|Pf-Qf|\le L\norm{P-Q}_{\TV}$.

\begin{lemma}[Predictable telescoping]\label{lem:telescoping}
Let $P$ and $Q$ be laws on the same finite visible history tree, with the
same initial law and the same feedback action kernels. Suppose their
conditional report kernels at history $h_t$ and action $a_t$ are
$p_t(\cdot\mid h_t,a_t)$ and $q_t(\cdot\mid h_t,a_t)$. If
\[
 \norm{p_t(\cdot\mid h_t,a_t)-q_t(\cdot\mid h_t,a_t)}_{\TV}
 \le d_t(h_t,a_t),
\]
then
\begin{equation}\label{eq:telescoping}
 \norm{P-Q}_{\TV}
 \le\min\left\{1,\E_P\sum_{t=0}^{N-1}d_t(H_t,A_t)\right\}.
\end{equation}
If both systems become absorbing under the same visible stopping rule
$\sigma$, the summand can be multiplied by $\1_{\{t<\sigma\}}$.
\end{lemma}
\begin{proof}
Fix a terminal measurable test $0\le f\le1$. Replace the report kernels
one at a time, using a $P$-prefix and a $Q$-suffix. The suffix conditional
expectation of $f$ is a measurable function between zero and one. The
change at step $t$, after integrating the common feedback action kernel,
is therefore bounded in absolute value by
$\E_P d_t(H_t,A_t)$. Summing the telescoping differences proves the bound
for $f$. Take the supremum over event indicators. After the visible rule
has stopped, the kernels are the same absorbing kernel and their local
error is zero. Histories of zero probability under $P$ contribute nothing;
the kernels are nevertheless specified on all admitted histories so that
all intermediate laws exist.
\end{proof}

The expectation is under the actual reference experiment. In particular a
rare report carrying a large conditional error is not automatically charged
as if it occurred at every time. No maximal-coupling measurability theorem is
needed for this proof. The action policy may be discontinuous: the same
kernel is applied on the same history in each telescoping term.

\begin{theorem}[Composition and risk transfer]\label{thm:composition}
Suppose $\Gamma_1:\cE\to\cF$ and $\Gamma_2:\cF\to\cG$ have
certificates $(\varepsilon_1,\rho_1)$ and
$(\varepsilon_2,\rho_2)$ uniformly over the controller classes needed by
the lifts. Their serial implementation has certificate
\[
 \bigl(\min\{1,\varepsilon_1+\varepsilon_2\},\,
       \rho_1\circ\rho_2\bigr).
\]
Exact reductions compose associatively as executions, up to relabeling of
product states. For a common target $\tau(\lambda)$ and bounded loss
$0\le\ell\le L$,
\begin{equation}\label{eq:risk-transfer}
 \mathfrak R_{\cE}(\rho_\Gamma(B))
 \le \mathfrak R_{\cF}(B)+L\varepsilon
\end{equation}
for either minimax risk or risk under a common prior.
\end{theorem}
\begin{proof}
Execute the second machine as a controller of the first. Their stored
states form a product and their source actions remain nonanticipative and
parameter independent. The first certificate applies to this enlarged
controller class by hypothesis. The subsequent machine/controller
execution is a Markov postprocessing of its visible history together with
its independent randomization, so it contracts total variation. Compare
first with an exact $\cF$ run and then with the intended $\cG$ run. The
triangle inequality gives the displayed defect. Budget maps compose in the
same order because they map target requirements to source requirements.
Regrouping three nested machine states changes only their names, proving
associativity.

For a target procedure within budget $B$, run its source implementation
and use the same final decision. Its risk at each parameter differs by at
most $L\varepsilon$. The resource certificate makes it admissible under
$\rho_\Gamma(B)$. Take the supremum over parameters, or integrate against
the common prior, and then take the infimum over target procedures. An
approximating sequence of procedures suffices if an optimum is not attained.
\end{proof}

A risk certificate concerns the same decision target and loss on both sides.
It does not compare excess errors relative to different full-history
predictors. An expected preparation constraint cannot be transported by a
pathwise certificate without checking its type. Conversely, total-variation
closeness alone does not transfer an unbounded cost expectation.

\subsection{Checking an exact state reduction locally}
\begin{proposition}[One-step realization]\label{prop:one-step}
Let $r_t:X_t\to X_t'$, $g_t:Y_t\to Y_t'$, and an action lift
$\iota_t:A_t'\to A_t$ be measurable. Suppose the preparations push forward
correctly and, for every admitted source state and target action,
\begin{equation}\label{eq:intertwining}
 (r_{t+1},g_{t+1})_\#
 K_t^\lambda(\cdot\mid x,\iota_ta')
 =K_t^{\prime\lambda}(\cdot\mid r_tx,a').
\end{equation}
Suppose also that continuation interfaces, recorded costs, and designated
readouts match. Then applying the lifted actions and report maps gives an
exact causal reduction for every target feedback strategy.
\end{proposition}
\begin{proof}
At time zero the reduced state laws agree. Conditional on a source state and
on a lifted action, \eqref{eq:intertwining} gives the joint reduced
next-state/report law. Integrate against the common current reduced law.
Since the target controller reads the reduced history, its next action law
is the same in both systems. Induction proves equality on all finite
cylinders, and hence on the path sigma field. Readout and cost matching give
the asserted experiment and resource identity.
\end{proof}

The state map need not be observable to implement the reduction: the
machine only needs the action lift and report map. The hidden-state map is
a proof device for the law identity. If the proposed lift instead requires
the unobserved state, it is not an executable reduction of this type.

\subsection{Terminal and causal equivalence}
For a fixed terminal family, a garbling $G$ is reversible if there exists a
parameter-independent kernel $R$ such that $P_\lambda GR=P_\lambda$.
This is statistical equivalence. Requiring $R$ to operate causally and
charging its memory gives a stronger notion. Requiring suitable bounds in
both directions gives resource-aware causal equivalence. These notions are
not identified by definition.

Independent parallel experiments may be implemented by product machines
when their preparations and physical randomness are genuinely independent
conditional on the parameter. Persistent cardinalities multiply and costs
add. Shared nuisance variables or coupled preparations require the joint
experiment; replacing them by an independent product changes the object.
